\makeatletter
\let\@oldaddcontentsline\addcontentsline
\renewcommand{\addcontentsline}[3]{}
\section*{附录}
\let\addcontentsline\@oldaddcontentsline
\makeatother

\subsection*{附录A：提交材料清单}

本论文附有以下支撑材料：

\begin{table}[H]
  \centering
  \caption{提交文件清单}
  \label{tab:appendix_files}
  \renewcommand{\arraystretch}{1.2}
  \begin{tabularx}{\textwidth}{LX}
    \toprule
    文件 & 说明 \\
    \midrule
    solve\_q1.m & 问题1 双光束模型推导与极值回归求解（v1.0）\\
    solve\_q2.m / run\_v15\_pipeline.m & 问题2 三域互证融合管线（v1.5，主程序）\\
    sic\_e1\_e2\_e3.m & 问题2 三域估计器核心实现 \\
    v20\_airy\_forward.m & 问题3 Airy 多光束前向模型（v2.0）\\
    v20\_tmm.m & 问题3 传输矩阵法验证实现 \\
    drude\_params.m & 问题3 Drude 介电函数参数 \\
    exp6\_w3\_si.m & 问题3 Si 判定实验（EXP-6）\\
    exp7\_sic\_w4\_exp8.m & 问题3 SiC 效应分离与精度量化（EXP-7/8）\\
    make\_paper\_figures.m & 论文图表生成脚本（MATLAB 基准图）\\
    draw\_all\_figures.py & 论文图表生成脚本（Python 重绘版）\\
    \bottomrule
  \end{tabularx}
\end{table}

\subsection*{附录B：核心代码（v1.0 双光束模型极值回归）}

以下为问题2 E1 极值回归法的核心 MATLAB 实现（完整管线代码另附）。代码以中文注释说明用途、输入输出与关键参数。

\begin{lstlisting}[language=Matlab, caption={双光束模型极值回归核心代码（E1）}]
% E1 极值回归法 — 条纹极值线性回归反演厚度
% 用途：从扣除背景后的条纹分量检测极值，线性回归得厚度
% 输入：sigma（波数列）, F（条纹分量列）, n1, theta_deg
% 输出：d_hat（厚度估计, um）, sigma_d（统计不确定度, um）
% 依赖：无

function [d_hat, sigma_d] = e1_extrema_regression(sigma, F, n1, theta_deg)

theta = deg2rad(theta_deg);
theta1 = asin(sin(theta)/n1);

% 亚像素极值检测：二次抛物线细化
[~, locs_max] = findpeaks(F);
[~, locs_min] = findpeaks(-F);
extrema_idx = sort([locs_max; locs_min]);

% 抛物线三点插值细化极值位置（亚像素）
for i = 1:length(extrema_idx)
    idx = extrema_idx(i);
    if idx > 1 && idx < length(sigma)
        xx = sigma(idx-1:idx+1);
        yy = F(idx-1:idx+1);
        p = polyfit(xx, yy, 2);
        sigma_refined(i) = -p(2)/(2*p(1));  % 抛物线顶点
    end
end

% 多假设周期竞争：在 GB/T 物理窗 [3,200] um 内取 top-5 候选
sigma_m = sigma_refined';
m = (1:length(sigma_m))';  % 起始级数赋值

% 线性回归：m = k*sigma_m + c
X = [sigma_m, ones(size(sigma_m))];
beta = X \ m;
k = beta(1);  % 斜率 = 4*n1*d*cos(theta1)
c = beta(2);  % 截距 = phi12/pi

d_hat = k / (4 * n1 * cos(theta1));  % 单位：cm
d_hat = d_hat * 1e4;  % 转换为 um

% 回归残差折算统计不确定度
res = m - X*beta;
sigma_d = std(res) / (4 * n1 * cos(theta1)) * 1e4;
end
\end{lstlisting}

\subsection*{附录C：Airy 多光束前向模型（v2.0）}

以下为问题3 的 Airy 多光束前向模型 MATLAB 核心实现。

\begin{lstlisting}[language=Matlab, caption={Airy 多光束前向模型核心代码（v2.0）}]
% v2.0 Airy 多光束前向模型 — 三介质 Airy 合成
% 用途：计算给定参数下的反射率谱 R(sigma)
% 输入：sigma（波数列, cm-1）, d（厚度, um）, n1（外延层折射率）,
%        theta_deg（入射角）, eps_inf, sigma_p, gamma（衬底 Drude 参数）
% 输出：R（反射率列, 0-1）
% 依赖：无

function R = v20_airy_forward(sigma, d, n1, theta_deg, ...
                               eps_inf, sigma_p, gamma)

theta = deg2rad(theta_deg);
theta1 = asin(sin(theta)/n1);

% 空气/外延层界面 Fresnel（标量）
r01 = (cos(theta) - n1*cos(theta1)) / (cos(theta) + n1*cos(theta1));
r10 = -r01;  % Stokes 关系

% 衬底 Drude 介电函数
eps_drude = eps_inf - sigma_p^2 ./ (sigma .* (sigma + 1i*gamma));
n2 = sqrt(eps_drude);  % 复折射率

% 衬底界面 Fresnel（复数）
theta2 = asin(n1*sin(theta1)./n2);
r12 = (n1*cos(theta1) - n2.*cos(theta2)) ./ ...
      (n1*cos(theta1) + n2.*cos(theta2));

% 往返相位
delta = 4*pi * n1 * (d*1e-4) * cos(theta1) * sigma;  % d 转 cm

% Airy 合成
E_r = r01 + (1-r01^2) * r12 .* exp(1i*delta) ./ ...
              (1 + r01 * r12 .* exp(1i*delta));

R = abs(E_r).^2;
% v2.0 退化验证：q=|r01*r12| -> 0 时精确退化为 v1.0 双光束式
% 数值偏差 3.9e-16（机器精度）
end
\end{lstlisting}